% =============================================================================
%  Section 4 -- Conclusions & Future Work
% =============================================================================

\section{Conclusions \& Future Work}

% --- Section divider --------------------------------------------------------
\sectiondivider{Conclusions \& Future Work}{What works, what limits us, and where this goes next}{sec:conclusions}

% --- Subsection: Hypothesis -------------------------------------------------
\subsection{Hypothesis Verdict}

\begin{frame}[label=sub:verdict]{Hypothesis verdict: \emph{partial support}}
\begin{exampleblock}{Evidence in favor}
\small
\begin{itemize}
    \item End-to-end pipeline works (video $\rightarrow$ pose $\rightarrow$ sequence $\rightarrow$ 5-class score).
    \item All three architectures land \alert{well above the 20\% random baseline}.
    \item Confusion matrices show ordinal errors -- the gradient is being learned.
    \item Multi-split sweep: 60/40 reaches \textbf{36.1\%}, $\rho\!=\!0.23$ ($p\!=\!0.001$).
\end{itemize}
\end{exampleblock}
\begin{alertblock}{Evidence against full validation}
\small
\begin{itemize}
    \item Documented final model (Test03): 29.9\% -- not deployment-ready.
    \item Generalization gap: $\sim$70~pp between training and test accuracy.
    \item Spearman $\rho$ at 80/20 not significant ($p\!>\!0.05$).
\end{itemize}
\end{alertblock}
{\footnotesize\itshape Proof of concept: the method works at the present scale; closing the remaining gap is a data \& labeling problem.}
\end{frame}

% --- Subsection: Limitations ------------------------------------------------
\subsection{Limitations}

\begin{frame}[label=sub:limitations]{Limitations}
\begin{itemize}\small
    \item \textbf{Dataset size}: 484 clips across 6 players -- minority-class supports as low as 8 and 13 in the test set.
    \item \textbf{Single-evaluator labels}: no inter-rater reliability computed (no Cohen's $\kappa$ or Intraclass Correlation Coefficient, ICC).
    \item \textbf{2D, 4 joints, 8 features}: no joint angles, no angular velocities, no lower-body, no depth.
    \item \textbf{Player leakage}: stratified split mixes the same player across train and test (no Leave-One-Player-Out, LOPO, evaluation yet).
    \item \textbf{Recording setup}: single lateral smartphone view $\Rightarrow$ self-occlusion + motion blur $\Rightarrow$ 11.6\% interpolation rate.
    \item \textbf{Categorical loss}: cross-entropy ignores the ordinal structure of the 5 quality classes.
\end{itemize}
\end{frame}

% --- Subsection: Future Work ------------------------------------------------
\subsection{Future Work \& Scalability}

\begin{frame}[label=sub:future]{Future work -- four highest-priority actions}
\begin{enumerate}\small
    \item \textbf{Dataset expansion} to $\geq$ 1{,}000 clips, target $\geq$ 200 per class, more players and varied environments.
    \item \textbf{Validated labeling protocol}: dual-grading on $\geq$ 20\% of clips, Cohen's $\kappa$ / ICC, consensus resolution.
    \item \textbf{Richer features}: joint angles, angular velocities, lower-body, 3D pose (multi-camera).
    \item \textbf{LOPO cross-validation} across players to bound true generalization.
\end{enumerate}
\vspace{6pt}
\textbf{\color{tecdark}Model-side directions:}
\begin{itemize}\small
    \item Adopt the \alert{XEnt + SGD} configuration identified in Sensitivity~1.
    \item Reduce augmentation factor; strengthen L2 / dropout (per Sensitivity~2).
    \item Try ordinal-aware losses, attention on top of BiLSTM~\cite{Vaswani2017Attention}, or Spatio-Temporal Graph Convolutional Networks (ST-GCN)~\cite{Yan2018STGCN}.
    \item Add phase-level diagnostics $\rightarrow$ first step toward a true recommendation module.
\end{itemize}
\end{frame}

\begin{frame}{Scalability of the pipeline}
\begin{columns}[T,onlytextwidth]
\begin{column}{0.47\textwidth}
\textbf{\color{tecdark}Other paddle tennis shots}
\begin{itemize}\small
    \item v\'ibora, smash, volley, globo.
    \item Different joints / sequence lengths / quality rubrics.
\end{itemize}
\vspace{8pt}
\textbf{\color{tecdark}Other racket sports}
\begin{itemize}\small
    \item Tennis, squash, pickleball, badminton.
    \item Same infrastructure; sport-specific data \& labels.
\end{itemize}
\end{column}
\hspace{0.02\textwidth}
\begin{column}{0.47\textwidth}
\textbf{\color{tecdark}Beyond racket sports}
\begin{itemize}\small
    \item Team sports (multi-person pose).
    \item Gymnastics, martial arts (codified scoring).
    \item Physical rehabilitation, ergonomics.
    \item Dance, performing arts.
\end{itemize}
\vspace{8pt}
\begin{infoBox}\scriptsize
The bottleneck is \alert{not} the architecture -- it is the quality and quantity of labeled data.
\end{infoBox}
\end{column}
\end{columns}
\end{frame}

% --- Subsection: Demo & Resources ------------------------------------------
\subsection{Demo \& Resources}

\begin{frame}[label=sub:demo]{Live demo \& open resources}
\begin{exampleblock}{Interactive Colab demo (MODEL\_KEY selector)}
\small \texttt{final} (Test03) -- \texttt{best\_grid} (XEnt + SGD) -- \texttt{best\_split} (60/40)\\[2pt]
\href{https://drive.google.com/file/d/1ilXqfHP970POqXvRvfTO6ZWcD54KAfLP/view?usp=drive_link}{\color{tecaccent}\textbf{$\rightarrow$ Open the demo on Google Drive}}
\end{exampleblock}
\vspace{6pt}
\textbf{\color{tecdark}Open-source code \& artifacts:}
\begin{itemize}\small
    \item Repo: \url{https://github.com/CesarEmilioC/thesisRepo_A00830006}
    \item Command-Line Interface (CLI): \texttt{pose}, \texttt{trainLSTM/GRU/TCN}, \texttt{predict*}, \texttt{plot}, \texttt{animate}, \texttt{runSensitivity}, \texttt{runSplitAnalysis}.
    \item Saved weights: \texttt{lstmModel\_Sens\_C\_xent\_sgd\_27-04-2026.h5}, \texttt{lstmModel\_Split60\_27-04-2026.h5}, plus Test03.
    \item Output artifacts: training history JSON, learning curves, confusion matrices, classification reports, class-distribution plots.
\end{itemize}
\end{frame}

% --- Subsection: Acronyms --------------------------------------------------
\subsection{Acronyms}

% ---- Acronyms page 1: Models, layers, training -----------------------------
\begin{frame}{Acronyms -- models, layers \& training}
\setlength{\extrarowheight}{1pt}
\setlength{\tabcolsep}{4pt}
\footnotesize
\begin{columns}[T,onlytextwidth]
\begin{column}{0.48\textwidth}
\textbf{\color{tecdark}Recurrent \& sequence models}\\[2pt]
\begin{tabular}{@{}l@{\hspace{8pt}}p{0.66\linewidth}@{}}
\textbf{LSTM}    & Long Short-Term Memory \\
\textbf{BiLSTM}  & Bidirectional LSTM \\
\textbf{GRU}     & Gated Recurrent Unit \\
\textbf{BiGRU}   & Bidirectional GRU \\
\textbf{RNN}     & Recurrent Neural Network \\
\textbf{TCN}     & Temporal Convolutional Network \\
\textbf{CNN}     & Convolutional Neural Network \\
\textbf{Conv1D}  & one-dimensional convolutional layer \\
\textbf{ST-GCN}  & Spatio-Temporal Graph Convolutional Network \\
\textbf{GAP}     & Global Average Pooling \\
\textbf{PAF}     & Part Affinity Field (in OpenPose) \\
\textbf{BPTT}    & Backpropagation Through Time \\
\end{tabular}
\end{column}
\hspace{0.02\textwidth}
\begin{column}{0.48\textwidth}
\textbf{\color{tecdark}Optimizers, losses \& regularization}\\[2pt]
\begin{tabular}{@{}l@{\hspace{8pt}}p{0.66\linewidth}@{}}
\textbf{Adam}    & Adaptive Moment estimation (optimizer) \\
\textbf{SGD}     & Stochastic Gradient Descent (optimizer) \\
\textbf{XEnt}    & Cross-entropy (loss function) \\
\textbf{MSE}     & Mean Squared Error (loss function) \\
\textbf{L2}      & L2 weight regularization (ridge / weight decay) \\
\end{tabular}\\[10pt]
\textbf{\color{tecdark}Quality classes (1--10 $\to$ 5)}\\[2pt]
\begin{tabular}{@{}l@{\hspace{8pt}}l@{}}
\textbf{1--2}    & Very Low \\
\textbf{3--4}    & Low \\
\textbf{5--6}    & Medium \\
\textbf{7--8}    & High \\
\textbf{9--10}   & Excellent \\
\end{tabular}
\end{column}
\end{columns}
\end{frame}

% ---- Acronyms page 2: Evaluation, statistics, tooling ----------------------
\begin{frame}{Acronyms -- evaluation, statistics \& tooling}
\setlength{\extrarowheight}{1pt}
\setlength{\tabcolsep}{4pt}
\footnotesize
\begin{columns}[T,onlytextwidth]
\begin{column}{0.48\textwidth}
\textbf{\color{tecdark}Evaluation \& statistics}\\[2pt]
\begin{tabular}{@{}l@{\hspace{8pt}}p{0.66\linewidth}@{}}
\textbf{TP}      & True Positive \\
\textbf{FP}      & False Positive \\
\textbf{FN}      & False Negative \\
\textbf{TN}      & True Negative \\
\textbf{F1}      & F1-score (harmonic mean of precision and recall) \\
\textbf{F1$_m$}  & Macro F1 (unweighted mean of per-class F1) \\
\textbf{$\rho$}  & Spearman's rank correlation coefficient \\
\textbf{$\kappa$}& Cohen's kappa (inter-rater agreement) \\
\textbf{ICC}     & Intraclass Correlation Coefficient \\
\textbf{pp}      & percentage points (absolute, not relative) \\
\textbf{LOPO}    & Leave-One-Player-Out (cross-validation) \\
\textbf{ROC}     & Receiver Operating Characteristic \\
\textbf{DET}     & Detection Error Tradeoff \\
\end{tabular}
\end{column}
\hspace{0.02\textwidth}
\begin{column}{0.48\textwidth}
\textbf{\color{tecdark}Hardware, data \& tooling}\\[2pt]
\begin{tabular}{@{}l@{\hspace{8pt}}p{0.66\linewidth}@{}}
\textbf{FPS}     & Frames Per Second \\
\textbf{IMU}     & Inertial Measurement Unit \\
\textbf{JSON}    & JavaScript Object Notation \\
\textbf{CLI}     & Command-Line Interface \\
\textbf{GUI}     & Graphical User Interface \\
\textbf{GPU}     & Graphics Processing Unit \\
\textbf{API}     & Application Programming Interface \\
\end{tabular}\\[10pt]
\textbf{\color{tecdark}Reading the metrics}\\[2pt]
{\scriptsize
A model with accuracy $a$ and Spearman $\rho$ near zero can rank classes randomly; one with modest $a$ but high $\rho$ preserves the quality ordering even when the exact class is wrong -- the regime we target on an ordinal problem.}
\end{column}
\end{columns}
\end{frame}

% --- Subsection: Glossary --------------------------------------------------
\subsection{Glossary}

% ---- Glossary page 1: Sport & vision (6 entries) ---------------------------
\begin{frame}{Glossary -- sport \& vision}
\footnotesize
\begin{description}\setlength{\itemsep}{8pt}\setlength{\parsep}{0pt}
\item[\textbf{Amateur}]\hfill\\
A sport player who participates primarily for enjoyment, personal satisfaction, and skill development rather than for financial gain.

\item[\textbf{Bandeja}]\hfill\\
Paddle tennis stroke typically executed in response to a mid-height lob from the opponent. The ideal hitting zone is mid-court, at a height between a volley and a smash.

\item[\textbf{Paddle tennis} (padel)]\hfill\\
Racket sport that merges aspects of tennis and squash. Played primarily in doubles on an enclosed court roughly one-third the size of a tennis court, with playable walls integral to the game.

\item[\textbf{Biomechanics}]\hfill\\
Study of the mechanical laws governing the movement and structure of living organisms; here, the analysis of joint trajectories, forces, and timing that characterize human motion in sports.

\item[\textbf{Pose estimation}]\hfill\\
Computer vision technique that detects and tracks human body key-points (e.g., joints) from images or videos to estimate spatial positions and orientations.

\item[\textbf{OpenPose}]\hfill\\
Open-source real-time multi-person 2D pose-estimation framework that predicts body, hand, facial, and foot key-points using confidence maps and Part Affinity Fields.
\end{description}
\end{frame}

% ---- Glossary page 2: Data & training-time concepts (5 entries) ------------
\begin{frame}{Glossary -- data \& training-time concepts}
\footnotesize
\begin{description}\setlength{\itemsep}{8pt}\setlength{\parsep}{0pt}
\item[\textbf{Stratified split}]\hfill\\
Partition of a dataset into training and test subsets that preserves the proportion of samples per class within each subset -- used here to keep the five quality classes balanced across both splits.

\item[\textbf{Data augmentation}]\hfill\\
Synthesizing additional training samples from existing ones by applying small perturbations (e.g., Gaussian noise on coordinate trajectories) to enlarge the effective training set and improve generalization on small datasets.

\item[\textbf{Cross-entropy loss}]\hfill\\
Standard classification loss that compares the predicted probability distribution against the one-hot encoded true label. Penalizes confident wrong predictions more strongly than uncertain ones.

\item[\textbf{Dropout}]\hfill\\
Regularization that randomly disables a fraction of neurons during each training step, forcing the network not to rely on any single feature path. This thesis uses 0.40 and 0.35 in the two BiLSTM layers.

\item[\textbf{L2 regularization} (weight decay)]\hfill\\
Penalty added to the loss proportional to the squared magnitudes of the model weights. Pushes weights toward zero, smooths decision boundaries, and helps when training data is scarce.
\end{description}
\end{frame}

% ---- Glossary page 3: Evaluation & generalization concepts (6 entries) -----
\begin{frame}{Glossary -- evaluation \& generalization}
\footnotesize
\begin{description}\setlength{\itemsep}{7pt}\setlength{\parsep}{0pt}
\item[\textbf{Ordinal classification}]\hfill\\
Classification in which target classes follow a natural ordering (\textit{Very Low} $<$ \textit{Low} $<$ \textit{Medium} $<$ \textit{High} $<$ \textit{Excellent}). Adjacent-class errors are less wrong than distant-class errors, motivating rank-based metrics such as Spearman's $\rho$.

\item[\textbf{Confusion matrix}]\hfill\\
$K\!\times\!K$ matrix where entry $(i, j)$ records true-class-$i$ samples predicted as class $j$. Diagonal entries are correct; off-diagonal entries reveal where errors land.

\item[\textbf{Overfitting}]\hfill\\
A model fitting the training data very closely -- including noise and player-specific idiosyncrasies -- but failing to generalize. Visible here as a $\sim$70 pp gap between train and test accuracy.

\item[\textbf{Inter-rater reliability}]\hfill\\
Statistical agreement between two or more annotators evaluating the same content (e.g., Cohen's $\kappa$, ICC). Not computed in this thesis -- each clip was graded by one coach.

\item[\textbf{Leave-one-player-out} (LOPO)]\hfill\\
Cross-validation protocol in which all clips of one player are held out for testing per fold; gives a conservative estimate of generalization to truly unseen individuals.

\item[\textbf{Percentage points} (pp)]\hfill\\
Absolute difference between two percentages: $26.8\%\!\to\!36.1\%$ is a gain of $9.3$~pp, not a $35\%$ relative increase.
\end{description}
\end{frame}

% --- Subsection: References -------------------------------------------------
\subsection{References}

\begin{frame}[allowframebreaks]{References}
\scriptsize
\setlength{\bibitemsep}{3pt}
\setlength{\bibhang}{8pt}
\printbibliography[heading=none]
\end{frame}

% --- Closing slide (Q&A) ----------------------------------------------------
\begin{frame}[plain]
\centering
\vspace{0.8cm}
{\color{tecblue}\rule{0.7\linewidth}{1.2pt}}\\[14pt]
{\Huge\bfseries\color{tecdark}Thank you}\\[10pt]
{\large\color{slidegray}\itshape Questions \& discussion}\\[14pt]
{\color{tecblue}\rule{0.7\linewidth}{1.2pt}}\\[20pt]
{\small Cesar Emilio Casta\~no Marin \quad A00830006}\\
{\scriptsize Tecnol\'ogico de Monterrey -- MCC -- May 2026}\\[8pt]
{\scriptsize\url{https://github.com/CesarEmilioC/thesisRepo_A00830006}}
\end{frame}
